\documentclass[12pt]{article}
\usepackage[utf8]{inputenc}
\usepackage[T1]{fontenc}
\usepackage{hyperref}
\title{Miahuatl\'an Zapotec Language Resources for K-12 Curriculum Development}
\author{Compiled by Claude Sonnet 5.5 (claude-sonnet-5-5) for Andr\'as Kornai}
\date{}
\begin{document}
\maketitle

\section{Introduction}
Miahuatl\'an Zapotec (ISO 639-3: zam) is an endangered Otomanguean language of Oaxaca, Mexico, spoken by an estimated 1,000 people. It is not known to be taught in schools. Little material exists, and this report says so plainly. Note that Mixtepec de Miahuatl\'an Zapotec (zpm) is a different language, and audio and scripture for it must not be mixed up with zam.

\section{Primary Educational Resources}
No primers, textbooks or school materials for zam were found in the searches.

\section{Linguistic and Reference Materials}
\begin{itemize}
\item Omniglot page on Miahuatl\'an Zapotec. It notes a dictionary, the Ruegsegger 1955 \emph{Vocabulario zapoteco del dialecto de Miahuatl\'an} published by the Instituto Ling\"u\'istico de Verano. I could not open the scanned copy it links to: \url{https://omniglot.com/writing/zapotec-miahuatlan.htm}
\item ASJP word-list entry: \url{https://asjp.clld.org/languages/MIAHUATLAN_ZAPOTEC}
\item Glosbe dictionary page (extent of vocabulary unverified): \url{https://glosbe.com/el/zam}
\end{itemize}

\section{Religious and Community Resources}
Omniglot says parts of the Bible have been translated. I did not confirm which portions exist or where they are hosted. A Joshua Project profile for the people group is at \url{https://joshuaproject.net/people_groups/11128}.

\section{Digital and Audio Resources}
Global Recordings Network lists evangelistic and Bible-teaching audio for zam: \url{https://globalrecordings.net/en/language/zam}

\section{Curriculum Development Context}
The language is endangered, and the only substantial published reference found is the 1950s SIL dictionary. A curriculum would need new documentation work with the community, starting from the dictionary and the audio recordings.

\end{document}